\subsection{The pairing of sieve observables with the Fourier boundary}
\label{acfp:section}

The logarithmic variable used by a sieve function records a divisor.
It does not record the prime counted by the resulting weight.
The incidence transform must therefore precede the passage to the prime measure.
The following construction gives that passage as an exact pairing.

\subsubsection{Removing the higher prime powers}

Fix a smooth function $\eta:\mathbb R\to[0,1]$ which vanishes for
$t\leq c$, with $0<c<\log2$, and equals one for $t\geq\log2$.
For a real parameter $\sigma$ and a compact smooth function $f$, define
\begin{equation}\label{acfp:mobius-test}
 (\mathcal P_\sigma f)(t)
   =\eta(t)\sum_{m\geq1}\mu(m)e^{-(m-1)\sigma t}f(mt).
\end{equation}
The summand is zero for all sufficiently large $m$, uniformly on $t\geq c$.
Consequently $\mathcal P_\sigma f$ is compact and smooth.
The choice of $\eta$ does not affect its pairing with a prime-power measure.

\begin{proposition}[A test map from prime powers to primes]\label{acfp:prime-projection}
For every compact smooth $f$ and every real $\sigma$,
\begin{equation}\label{acfp:projection-pairing}
 \mu_\sigma(\mathcal P_\sigma f)=\nu_\sigma(f).
\end{equation}
If $s=\sigma+a\geq1$, then $\mathcal P_\sigma$ extends to a continuous
linear map $\mathcal S_a\to\mathcal S_a$ and the same identity holds there.
Under the damping isomorphism,
\[
 M_a\mathcal P_\sigma M_a^{-1}=\mathcal Q_s,\qquad
 (\mathcal Q_sg)(t)
    =\eta(t)\sum_{m\geq1}\mu(m)e^{-(m-1)st}g(mt).
\]
For every pair of nonnegative integers $r,j$, a constant $C_{r,j,s,\eta}$
satisfies
\begin{equation}\label{acfp:projection-norm}
 \sup_t(1+|t|)^r|\partial_t^j\mathcal Q_sg(t)|
 \leq C_{r,j,s,\eta}
       \sum_{\ell=0}^j\sup_u(1+|u|)^r|g^{(\ell)}(u)|.
\end{equation}
\end{proposition}

\begin{proof}
The compact pairing is a finite sum.  Expanding it gives
\[
 \sum_{p,j,m\geq1}
   (\log p)\mu(m)p^{-jm\sigma}f(jm\log p)
 =\sum_{p,k\geq1}(\log p)p^{-k\sigma}f(k\log p)
                      \sum_{m\mid k}\mu(m).
\]
The inner divisor sum is one for $k=1$ and zero otherwise.
This proves \eqref{acfp:projection-pairing}.
No primality test enters the operator $\mathcal P_\sigma$.

Conjugation by $M_a$ gives the stated formula for $\mathcal Q_s$.
For $t\geq c$, the product rule bounds each derivative of a summand by
a finite sum of
\[
 C_{j,s,\eta}m^j e^{-(m-1)sc}|g^{(\ell)}(mt)|,\qquad 0\leq\ell\leq j.
\]
All derivatives of $\eta$ are bounded.
Also $(1+t)^r\leq(1+mt)^r$ there.
The convergent series $\sum_{m\geq1}m^j e^{-(m-1)sc}$ proves
\eqref{acfp:projection-norm} and convergence in every Schwartz seminorm.
On compact tests, $\nu_\sigma(f)=\nu_s(M_af)$.
Theorem~\ref{acfd:tempered} therefore extends $\nu_\sigma$ continuously
to $\mathcal S_a$.
Theorem~\ref{acfd:exponential-domain} and the bound for $\mathcal Q_s$
give continuity of the other side of \eqref{acfp:projection-pairing}.
Density of compact smooth functions extends the identity.
\end{proof}

\subsubsection{The exact observable map}

Use the finite arithmetic parameters in Theorem~\ref{acfd:exact} and set
\[
 \rho=\frac{\log R}{\log N}\in(0,1/2),\qquad
 C_N=\frac{\varphi(W)^kN(\log R)^k}{W^{k+1}}.
\]
Write $\mathcal D_N,\Lambda_N,T_N,I_N$ for the corresponding finite objects.
In the asymptotic regime, take the previously specified $R=R_N=N^\rho$
and $W=W_N$, with $\rho>0$ fixed.
Let $\chi\in C_c^\infty((-1,1);\mathbb R)$ satisfy
$0\leq\chi\leq1$ and $\chi(0)=1$.
This $\chi$ denotes the interpolation bump in this subsection.
Set $H=\max_i h_i$ and $r_N=[4(2N+H+1)]^{-1}$.

For $\lambda,\gamma\in\Lambda_N$ and a real observable
$w:I_N\to\mathbb R$, define a compact smooth function
\begin{equation}\label{acfp:observable-map}
 \begin{split}
 \mathcal B_{N,\sigma}(\lambda,\gamma;w)(t)
  ={}&\frac1{\rho C_N}
    \sum_{i=1}^k\sum_{n\in I_N}
       w(n)(T_N\lambda)(n)(T_N\gamma)(n)
       \frac{(n+h_i)^\sigma}{\log(n+h_i)}\\
     &\hspace{35mm}\cdot
       \chi\left(\frac{t-\log(n+h_i)}{r_N}\right).
 \end{split}
\end{equation}
The map is bilinear in $(\lambda,\gamma)$ and linear in $w$.
It contains no prime indicator.
The domain is the original coefficient space and the original counting space.

\begin{theorem}[An exact pairing with the prime boundary]\label{acfp:exact-observable}
Let $\sigma,a\geq0$ satisfy $\sigma+a=1$.
Set
\[
 \Theta_N(\lambda,\gamma;w)
   =\mathcal F^{-1}M_a\mathcal P_\sigma
                      \mathcal B_{N,\sigma}(\lambda,\gamma;w)
       \in\mathcal S(\mathbb R).
\]
Then
\begin{equation}\label{acfp:spine}
 \widehat\mu_1\bigl(\Theta_N(\lambda,\gamma;w)\bigr)
   =\frac1{\rho C_N}\sum_{n\in I_N}
       w(n)R_{\mathcal H}(n)(T_N\lambda)(n)(T_N\gamma)(n).
\end{equation}
In particular, $\lambda=E_NF$, $\gamma=E_NG$, and $w=1$ give
the exact prime form $b_N(F,G)$.

The support of $\mathcal B_{N,\sigma}$ lies in
$[\log N-r_N,\log(2N+H)+r_N]$.
For $\sigma\geq0$ and each derivative order $j$,
\begin{equation}\label{acfp:observable-norm}
 \|\partial_t^j\mathcal B_{N,\sigma}(\lambda,\gamma;w)\|_\infty
 \leq
 \frac{k\|w\|_\infty\|\lambda\|_1\|\gamma\|_1}
      {\rho C_N}
 \frac{(2N+H)^\sigma}{\log N}\,r_N^{-j}\|\chi^{(j)}\|_\infty.
\end{equation}
\end{theorem}

\begin{proof}
Each bump meets just one logarithm of an integer, by the spacing estimate
in Proposition~\ref{acfd:log-realization}.
Its pairing with $\nu_\sigma$ is one if its center is a prime and zero otherwise.
The factor $(n+h_i)^\sigma/\log(n+h_i)$ cancels the mass at that prime.
Thus the right side of \eqref{acfp:spine} equals
$\nu_\sigma(\mathcal B_{N,\sigma})$.
Proposition~\ref{acfp:prime-projection} makes this
$\mu_\sigma(\mathcal P_\sigma\mathcal B_{N,\sigma})$.
Since $\sigma+a=1$, Theorem~\ref{acfd:exponential-domain} makes it
$\mu_1(M_a\mathcal P_\sigma\mathcal B_{N,\sigma})$.
The Fourier convention then gives the left side exactly, including the
factor $(2\pi)^{-1}$ in $\mathcal F^{-1}$.

At any center, at most $k$ pairs $(i,n)$ occur.
Distinct integer centers have disjoint bump supports.
Also $|T_N\lambda(n)|\leq\|\lambda\|_1$.
These facts prove \eqref{acfp:observable-norm} and the support statement.
\end{proof}

\begin{corollary}[The coefficient quotient]\label{acfp:quotient}
The test maps $\mathcal B_{N,\sigma}$ and $\Theta_N$ descend to
$\mathcal Q_N\times\mathcal Q_N\times\mathbb R^{I_N}$, where
$\mathcal Q_N=\Lambda_N/\ker T_N$ has norm
$\|[\lambda]\|_N=C_N^{-1/2}\|T_N\lambda\|_2$.
They are continuous for each fixed $N$.
In particular,
\[
 \|\partial_t^j\mathcal B_{N,\sigma}(\lambda,\gamma;w)\|_\infty
 \leq\frac{k\|w\|_\infty}{\rho}
 \|[\lambda]\|_N\|[\gamma]\|_N
 \frac{(2N+H)^\sigma}{\log N}\,r_N^{-j}\|\chi^{(j)}\|_\infty.
\]
\end{corollary}

\begin{proof}
The formula uses a coefficient vector only through its incidence values.
It therefore has the asserted quotient domain.
At each $n$, the product of incidence values is bounded by
$\|T_N\lambda\|_2\|T_N\gamma\|_2$.
Substitute this bound in the proof of \eqref{acfp:observable-norm}.
The common compact support bounds every weighted derivative seminorm.
The continuous maps $M_a$, $\mathcal P_\sigma$, and $\mathcal F^{-1}$
give continuity of $\Theta_N$.
\end{proof}

This quotient map has a distributional test space as codomain.
It does not construct the map $R_N^\sharp:\mathcal Q_N\to H_k$
or the uniform quadratic inequality in Proposition~\ref{prop:full-criterion}.

For $F,G\in\mathcal X_k$, the composite in Theorem~\ref{acfp:exact-observable}
is therefore a specified bilinear map from the smooth simplex carrier
to the Schwartz test space of the prime-power Fourier boundary.
Under \eqref{eq:lod}, Theorem~\ref{thm:forms} gives
\begin{equation}\label{acfp:conditional-limit}
 \widehat\mu_1\bigl(\Theta_N(E_NF,E_NG;1)\bigr)
       \longrightarrow\langle F,K_kG\rangle.
\end{equation}
This is convergence of pairings with a varying family of tests.
It does not assert convergence of those tests in Schwartz topology.

\subsubsection{A finite first moment on the counting space}

Take $k=1$, $\mathcal H=\{0\}$, $N=256$, $R=4$, $W=2$, and $v=1$.
Then $\mathcal D_N=\{1,3\}$ and $\rho=1/4$.
Choose $F\in C_c^\infty((0,1))$ with
$F(\log3/\log4)=1$.
A smooth bump about this interior point gives such a function.
The actual sample vector is $P_NF=(0,1)$, so
\[
 L_N=\begin{pmatrix}1&1/2\\0&-3/2\end{pmatrix},\qquad
 \lambda=E_NF=(1/2,-3/2),\qquad
 (T_N\lambda)(n)=\frac12-\frac32\mathbf1_{3\mid n}.
\]
Among the $128$ odd integers in $[256,512)$, exactly $42$ are divisible by $3$.
The interval contains $43$ primes, whose sum is $16467$.
Trial division by $2,3,5,7,11,13,17,19$ determines these counts,
since $23^2>512$.
For $w_0(n)=1$ and $w_1(n)=n/N$, finite summation gives
\[
\begin{array}{c|cc}
 & w_0 & w_1\\\hline
 \displaystyle\sum_{n\in I_N}w_j(n)(T_N\lambda)(n)^2
       &127/2&381/4\\[1mm]
 \displaystyle\sum_{p\in I_N\cap\mathbb P}w_j(p)(T_N\lambda)(p)^2
       &43/4&16467/1024
\end{array}
\]
Here $C_N=128\log2$ and $\rho C_N=32\log2$.
Thus the exact boundary pairings at $\sigma=a=1/2$ are
\begin{equation}\label{acfp:finite-moment}
 \widehat\mu_1(\Theta_N(\lambda,\lambda;w_0))
       =\frac{43}{128\log2},\qquad
 \widehat\mu_1(\Theta_N(\lambda,\lambda;w_1))
       =\frac{16467}{32768\log2}.
\end{equation}
The normalized first moment of the prime-weighted counting measure is
\[
 \frac{16467}{11008}=\frac32-\frac{45}{11008}.
\]
Its mass-weighted counterpart is $3/2$.
The pairing preserves the exact arithmetic moment, including this difference.

\subsubsection{A nonconstant one-dimensional model}

Take $k=1$, $\mathcal H=\{0\}$, and the polynomial $f(t)=t$ on $[0,1]$.
Then $T_1f=\int_0^1t\,dt=1/2$, so
\begin{equation}\label{acfp:linear-model}
 I_f=\int_0^1t^2\,dt=\frac13,\qquad
 \langle f,K_1f\rangle=\left(\int_0^1t\,dt\right)^2=\frac14,
 \qquad R_1(f)=\frac34.
\end{equation}
Let
$\varepsilon_N=\min\{1/8,D_0(N)^{-2/5}\}>0$,
$F_N=\Xi_{\varepsilon_N}f$, and $\lambda_N=E_NF_N$.
These are the positive-width functions in Corollary~\ref{acfd:native-polynomial}.
Assume the same distribution hypothesis and fixed $0<\rho<\vartheta/2$.
Then the actual finite forms satisfy
\begin{equation}\label{acfp:linear-limits}
 \frac{S_1(\lambda_N)}{C_N}\longrightarrow\frac13,\qquad
 \widehat\mu_1\bigl(\Theta_N(\lambda_N,\lambda_N;1)\bigr)
       =\frac{S_2(\lambda_N)}{\rho C_N}\longrightarrow\frac14,
 \qquad
 \frac{S_2(\lambda_N)}{S_1(\lambda_N)}\longrightarrow\frac{3\rho}4.
\end{equation}
The two form limits follow from the quantitative estimates in the proof
of Theorem~\ref{acfd:moving-cutoff}, applied before dividing the forms.

The constant and linear simplex functions give the full first comparison.
Put $f_0(t)=1$, $f_1(t)=t$, and
$\lambda_{j,N}=E_N(\Xi_{\varepsilon_N}f_j)$.
The two limiting matrices, indexed by $i,j\in\{0,1\}$, are
\begin{equation}\label{acfp:first-matrices}
 \left(\frac{S_1(\lambda_{i,N},\lambda_{j,N})}{C_N}\right)_{i,j}
 \longrightarrow
 \begin{pmatrix}1&1/2\\1/2&1/3\end{pmatrix},\qquad
 \left(\widehat\mu_1(\Theta_N(\lambda_{i,N},\lambda_{j,N};1))\right)_{i,j}
 \longrightarrow
 \begin{pmatrix}1&1/2\\1/2&1/4\end{pmatrix}.
\end{equation}
Indeed, the bilinear error bounds in Theorem~\ref{thm:forms} are
$O(D_0^{-1}\varepsilon_N^{-2})$ for this fixed pair of polynomials.
The cutoff errors tend to zero in $L^2$.
The limiting entries are
$\int_0^1 f_if_j$ and $(\int_0^1f_i)(\int_0^1f_j)$, respectively.
The mass matrix has determinant $1/12$, whereas the prime matrix has rank one.
Thus the boundary pairing matches the fiber form, with its different rank
and second entry, on these original simplex functions.

There is also an exact arithmetic simplification.
For every prime $p\in I_N$, a divisor $d<R_N<N$ of $p$ must equal one.
Hence
\begin{equation}\label{acfp:one-coordinate}
 (T_N\lambda_N)(p)=(\lambda_N)_1,\qquad
 S_2(\lambda_N)
    =\#(I_N\cap\mathbb P)\left|
       \sum_{r\in\mathcal D_N}\frac{F_N(\log r/\log R_N)}{\varphi(r)}
                         \right|^2.
\end{equation}
The divisor average is taken before its square weights the primes.

\begin{proposition}[The direct prime-variable substitution fails]\label{acfp:substitution}
For the same parameters, the expression obtained by substituting the
prime variable directly into the simplex function is identically zero:
\[
 \frac1{\rho C_N}\sum_{p\in I_N\cap\mathbb P}
                |F_N(\log p/\log R_N)|^2=0.
\]
It therefore differs from the actual normalized prime form in
\eqref{acfp:linear-limits}, whose conditional limit is $1/4$.
\end{proposition}

\begin{proof}
Every prime in $I_N$ is at least $N$.
Since $R_N=N^\rho$ and $\rho<1/2$,
$\log p/\log R_N\geq1/\rho>1$.
The zero extension of $F_N$ vanishes there.
The actual form first applies $L_NP_N$ and then incidence, as in
\eqref{acfp:one-coordinate}.  Its limit is \eqref{acfp:linear-limits}.
\end{proof}

\subsubsection{Why the exact tests have no bounded Schwartz limit}

The failure of direct substitution does not invalidate the exact pairing.
Its moving tests have a quantitative topology obstruction.
In the preceding model, take $\sigma=a=1/2$, and put
\[
 g_N=M_{1/2}\mathcal B_{N,1/2}(\lambda_N,\lambda_N;1),\qquad
 u_N=\mathcal Q_1g_N,\qquad \Theta_N=\mathcal F^{-1}u_N.
\]
These functions are Schwartz.  The function $g_N$ is nonnegative and has
support in $J_N=[\log N-r_N,\log(2N)+r_N]$.
The notation $J_N$ in this paragraph denotes this real interval only.

\begin{proposition}[A growing seminorm]\label{acfp:unbounded-tests}
Under the hypotheses of \eqref{acfp:linear-limits}, there is $c_0>0$
such that, for all sufficiently large $N$,
\[
 \sup_t(1+|t|)^2|g_N(t)|\geq c_0\log N,\qquad
 \sup_t(1+|t|)^2|u_N(t)|\geq c_0\log N.
\]
Consequently neither $(u_N)$ nor $(\Theta_N)$ is bounded in Schwartz space.
\end{proposition}

\begin{proof}
The exact pairing and damping identity give
$\nu_1(g_N)=S_2(\lambda_N)/(\rho C_N)\geq1/8$ eventually.
The interval $J_N$ has length less than two for large $N$.
It meets at most three intervals $[j,j+1)$ with integer $j\geq0$.
The proof of Theorem~\ref{acfd:tempered} therefore gives
\[
 \nu_1(J_N)\leq\mu_1(J_N)
       \leq3e(\log(2N)+3).
\]
Positivity now implies
$\|g_N\|_\infty\geq[24e(\log(2N)+3)]^{-1}$.
Since every point of its support is at least $\log N-r_N$, its second
weighted seminorm grows at least as a positive constant times $\log N$.

For $t\in J_N$ and large $N$, one has $t\geq\log2$ and
$2t>\log(2N)+r_N$.
Thus $\eta(t)=1$ and $g_N(mt)=0$ for every $m\geq2$.
Formula \eqref{acfp:mobius-test}, after damping, gives
$u_N(t)=g_N(t)$ throughout $J_N$.
The same lower bound applies to $u_N$.
Fourier transformation is a topological automorphism of Schwartz space,
so boundedness of $(\Theta_N)$ would imply boundedness of $(u_N)$.
\end{proof}

\subsubsection{A controlled approach through the holomorphic half-plane}

An unbounded family of tests requires a separate estimate when the
distribution also varies.
For the finite data in Theorem~\ref{acfp:exact-observable}, put
$L_N^{\mathrm{log}}=\log(2N+H)$ and define the nonnegative number
\begin{equation}\label{acfp:boundary-control}
 U_N=
 \frac{k\|w\|_\infty\|\lambda\|_1\|\gamma\|_1}
      {\rho C_N(1-e^{-c})}
 \frac{(2N+H)e^{ar_N}}{\log N}
       (1+L_N^{\mathrm{log}}+r_N)^4.
\end{equation}
The constant $c>0$ is the cutoff parameter in
\eqref{acfp:mobius-test}.

\begin{proposition}[A diagonal logarithmic-derivative pairing]
\label{acfp:holomorphic-pairing}
For every $\epsilon>0$, the following ordinary integral converges absolutely:
\[
 Z_{N,\epsilon}
   =\int_{\mathbb R}
       -\frac{\zeta'(1+\epsilon+i\xi)}{\zeta(1+\epsilon+i\xi)}
          \Theta_N(\lambda,\gamma;w)(\xi)\,d\xi.
\]
It satisfies
\begin{equation}\label{acfp:holomorphic-error}
 \left|Z_{N,\epsilon}
   -\frac1{\rho C_N}\sum_{n\in I_N}
      w(n)R_{\mathcal H}(n)(T_N\lambda)(n)(T_N\gamma)(n)\right|
       \leq C_*\epsilon U_N.
\end{equation}
For any sequence of the finite data, the explicit choice
$\epsilon_N=[N(1+U_N)]^{-1}>0$ makes this error at most $C_*/N$.
For the linear model \eqref{acfp:linear-limits}, this choice gives
$Z_{N,\epsilon_N}\to1/4$ under its stated distribution hypothesis.
\end{proposition}

\begin{proof}
Write $g=M_a\mathcal B_{N,\sigma}(\lambda,\gamma;w)$ and
$u=\mathcal Q_1g$.
Then $\widehat\Theta_N=u$.
The bounded function $G_{1+\epsilon}$ in
Theorem~\ref{acfd:fourier-boundary} represents
$\widehat\mu_{1+\epsilon}$, and $\Theta_N$ is integrable.
Hence the ordinary integral equals
$\widehat\mu_{1+\epsilon}(\Theta_N)$.

For $t\geq c$ and $m\geq1$, one has
$(1+t)^4\leq(1+mt)^4$ and
$e^{-(m-1)t}\leq e^{-(m-1)c}$.
Since $0\leq\eta\leq1$, the geometric series gives
\[
 \sup_t(1+|t|)^4|u(t)|
       \leq\frac1{1-e^{-c}}\sup_t(1+|t|)^4|g(t)|.
\]
On the bump centered at $\log(n+h_i)$, the damping factor is at most
$(n+h_i)^a e^{ar_N}$.
Use $\sigma+a=1$, the disjoint supports, and at most $k$ equal centers.
The resulting fourth seminorm is bounded by $U_N$.
Equation~\eqref{acfd:boundary-limit} now gives
$|Z_{N,\epsilon}-\widehat\mu_1(\Theta_N)|\leq C_*\epsilon U_N$.
The exact pairing proves \eqref{acfp:holomorphic-error}.
The selected $\epsilon_N$ tends to zero and satisfies
$\epsilon_NU_N\leq1/N$.
Finally apply \eqref{acfp:linear-limits}.
\end{proof}

Equations~\eqref{acfp:conditional-limit} and \eqref{acfp:holomorphic-error}
compare the simplex fiber form with the prime boundary through incidence
and a controlled diagonal limit.
A complete explicit formula still requires its other terms and convergence
estimates on a common test domain.
